\documentclass[11pt,a4paper]{article}

% Codificación, idioma y tipografía
\usepackage[utf8]{inputenc}
\usepackage[T1]{fontenc}
\usepackage[spanish,es-noquoting]{babel}
\usepackage{lmodern}
\usepackage{microtype}

% Matemáticas, tablas y diseño
\usepackage{amsmath,amssymb}
\usepackage{geometry}
\usepackage{booktabs,tabularx,array}
\usepackage{enumitem}
\usepackage{xcolor}
\usepackage{tikz}
\usetikzlibrary{arrows.meta,positioning}
\usepackage{hyperref}
\geometry{margin=2.5cm}
\hypersetup{colorlinks=true,linkcolor=blue!60!black,urlcolor=blue!60!black}
\setlist[itemize]{topsep=3pt,itemsep=3pt}
\setlength{\parskip}{5pt}
\setlength{\parindent}{0pt}
\renewcommand{\arraystretch}{1.25}

\title{Fundamentos del algoritmo de Grover\\[2mm]
\large Evaluación bajo distintos modelos de ruido}
\author{Trabajo de Fin de Grado --- Documento de estudio}
\date{\today}

\begin{document}
\maketitle
\tableofcontents
\bigskip

\section{Objetivo de esta fase}

El objetivo de este primer apartado es comprender \textbf{cómo funciona el algoritmo de Grover}, por qué es importante elegir correctamente el número de iteraciones y cómo puede afectar el ruido a su comportamiento. También se identifican las herramientas necesarias para su futura implementación.

En esta fase se estudian los fundamentos teóricos; todavía no es necesario programar los circuitos.

\section{¿Qué es el algoritmo de Grover?}

El algoritmo de Grover, propuesto por Lov Grover en 1996, es un algoritmo cuántico de \textbf{búsqueda no estructurada}. Su principal ventaja consiste en reducir el número de consultas al oráculo necesarias para encontrar una solución.

\begin{center}
\begin{tabular}{ll}
\toprule
\textbf{Método} & \textbf{Complejidad de consultas} \\
\midrule
Búsqueda clásica no estructurada & $O(N)$ \\
Algoritmo de Grover & $O(\sqrt{N})$ \\
\bottomrule
\end{tabular}
\end{center}

Aquí, $N$ es el número de elementos del espacio de búsqueda. Por ejemplo, si $N=1\,000\,000$, una búsqueda clásica puede requerir del orden de un millón de comprobaciones, mientras que Grover necesita del orden de un millar de consultas al oráculo.

\textbf{Importante:} esta ventaja corresponde al \emph{número de consultas al oráculo}; no implica que la emulación de Grover en un ordenador clásico sea más rápida que una búsqueda clásica equivalente.

\section{Conceptos cuánticos necesarios}

\begin{center}
\begin{tabularx}{\textwidth}{@{}>{\bfseries}p{3.0cm}X@{}}
\toprule
Concepto & Explicación \\
\midrule
Qubit & Unidad de información cuántica que puede estar en superposición de los estados $\lvert 0\rangle$ y $\lvert 1\rangle$. \\
Superposición & Combinación de estados básicos mediante amplitudes cuánticas. \\
Puertas cuánticas & Operaciones que transforman el estado de uno o varios qubits. \\
Interferencia & Fenómeno que permite reforzar o cancelar amplitudes. \\
Medición & Proceso que produce un resultado clásico según las probabilidades del estado cuántico. \\
\bottomrule
\end{tabularx}
\end{center}

Las puertas más relevantes para el proyecto son:

\begin{itemize}
    \item \textbf{H (Hadamard):} permite crear superposiciones.
    \item \textbf{X (Pauli-X):} intercambia $\lvert 0\rangle$ y $\lvert 1\rangle$.
    \item \textbf{Z (Pauli-Z):} modifica la fase de $\lvert 1\rangle$.
    \item \textbf{CZ y puertas controladas o multicontroladas:} facilitan la construcción de oráculos y operadores de difusión.
\end{itemize}

No es necesario dominar toda la computación cuántica para empezar, pero sí entender cómo se utilizan estas operaciones en el algoritmo.

\section{Funcionamiento del algoritmo de Grover}

El procedimiento comprende cuatro etapas:

\begin{center}
\begin{tikzpicture}[
  node distance=8mm,
  block/.style={draw=blue!55!black,fill=blue!4,rounded corners=3pt,
                align=center,text width=10.8cm,minimum height=1.0cm},
  arr/.style={-{Latex[length=2.6mm]},thick,draw=blue!65!black}
]
\node[block] (ini) {\textbf{1. Inicialización}\\Superposición uniforme de los estados posibles};
\node[block,below=of ini] (ora) {\textbf{2. Oráculo}\\Cambio de fase del estado que cumple la condición};
\node[block,below=of ora] (dif) {\textbf{3. Difusión}\\Amplificación de la amplitud de la solución};
\node[block,below=of dif] (med) {\textbf{4. Medición}\\Obtención de un resultado clásico};
\draw[arr] (ini) -- (ora);
\draw[arr] (ora) -- (dif);
\draw[arr] (dif) -- (med);
\draw[arr] (dif.east) -- ++(9mm,0) |- node[right,pos=0.30,font=\footnotesize]{Repetir iteraciones} (ora.east);
\end{tikzpicture}
\end{center}

\subsection{Ejemplo con dos qubits}

Con dos qubits se dispone de cuatro estados posibles:
\[
\bigl\{\lvert 00\rangle,\lvert 01\rangle,\lvert 10\rangle,\lvert 11\rangle\bigr\}.
\]

Supongamos que el objetivo es encontrar el estado $\lvert 11\rangle$.

\textbf{Paso 1. Inicialización.} Aplicando una puerta Hadamard a cada qubit, se obtiene
\[
\lvert\psi\rangle=\frac12\bigl(\lvert 00\rangle+\lvert 01\rangle+\lvert 10\rangle+\lvert 11\rangle\bigr).
\]
Cada estado tiene inicialmente una probabilidad del $25\%$ de ser medido.

\textbf{Paso 2. Oráculo.} Se cambia el signo de la amplitud asociada al estado buscado:
\[
\lvert\psi'\rangle=\frac12\bigl(\lvert 00\rangle+\lvert 01\rangle+\lvert 10\rangle-\lvert 11\rangle\bigr).
\]
El cambio de fase no modifica por sí solo las probabilidades de medición.

\textbf{Paso 3. Difusión.} Se aplica una reflexión de las amplitudes alrededor de su media. En este ejemplo, las amplitudes de los estados no deseados pasan a cero y la de $\lvert 11\rangle$ pasa a uno.

\textbf{Paso 4. Medición.} En ausencia de ruido, una iteración permite obtener $\lvert 11\rangle$ con probabilidad del $100\%$.

\begin{center}
\begin{tabular}{@{}lcc@{}}
\toprule
\textbf{Estado} & \textbf{Antes de Grover} & \textbf{Tras una iteración} \\
\midrule
$\lvert 00\rangle$ & $25\%$ & $0\%$ \\
$\lvert 01\rangle$ & $25\%$ & $0\%$ \\
$\lvert 10\rangle$ & $25\%$ & $0\%$ \\
$\lvert 11\rangle$ & $25\%$ & $100\%$ \\
\bottomrule
\end{tabular}
\end{center}

\section{Número de iteraciones de Grover}

Una \emph{iteración de Grover} aplica el oráculo seguido del operador de difusión. El número de iteraciones adecuado depende del tamaño del espacio de búsqueda.

Para \textbf{una única solución marcada}, la aproximación habitual es
\[
\boxed{r_{\mathrm{opt}}\approx\frac{\pi}{4}\sqrt{N}.}
\]

La probabilidad de éxito ideal después de $r$ iteraciones viene dada por
\[
\boxed{P(r)=\sin^2\!\bigl((2r+1)\theta\bigr),\qquad
\theta=\arcsin\!\left(\frac{1}{\sqrt{N}}\right).}
\]

Para un análisis más preciso, se selecciona un número entero de iteraciones cercano a
\[
r\approx \frac{\pi}{4\theta}-\frac12.
\]

\subsection{Ejemplo con tres qubits}

Con tres qubits hay $N=2^3=8$ estados y, si solo existe una solución, la probabilidad inicial de acertar es del $12{,}5\%$.

\begin{center}
\begin{tabular}{@{}cc@{}}
\toprule
\textbf{Iteraciones ($r$)} & \textbf{Probabilidad ideal de éxito} \\
\midrule
0 & $12{,}50\%$ \\
1 & $78{,}13\%$ \\
2 & $94{,}53\%$ \\
3 & $33{,}01\%$ \\
4 & $1{,}22\%$ \\
5 & $54{,}80\%$ \\
6 & $99{,}98\%$ \\
\bottomrule
\end{tabular}
\end{center}

Dos iteraciones ofrecen un valor muy elevado dentro del \emph{primer máximo} de amplificación. La probabilidad \textbf{no aumenta de manera monótona}: si se siguen aplicando iteraciones, disminuye y posteriormente puede aumentar de nuevo. Por tanto, realizar más iteraciones no garantiza una mejora.

\section{Modelos de ruido}

En un circuito ideal, las puertas y las mediciones se consideran perfectas. En cambio, los sistemas cuánticos reales están expuestos a errores. En este TFG se reproducirán mediante \textbf{modelos de ruido simulados}.

\begin{center}
\begin{tabularx}{\textwidth}{@{}>{\bfseries}p{3.0cm}X X@{}}
\toprule
Modelo & Qué representa & Posible efecto sobre Grover \\
\midrule
Despolarización & Alteraciones del estado asociadas a operaciones imperfectas. & Deterioro de la amplificación y acumulación de errores cuando aumenta el número de puertas. \\
Error de lectura & Errores al registrar los bits obtenidos en la medición. & Observación de un estado incorrecto aunque el circuito haya preparado el correcto. \\
\bottomrule
\end{tabularx}
\end{center}

El \textbf{ruido despolarizante} aplicado a puertas resulta especialmente interesante para estudiar las iteraciones: aumentar su número añade operaciones y, con ello, más oportunidades de sufrir errores.

Los \textbf{errores de lectura}, si solo se introducen al final del circuito, no se acumulan durante las iteraciones. Por ello, conviene analizarlos por separado.

\subsection{Hipótesis del TFG}

\begin{quote}
\emph{Bajo determinadas condiciones de ruido, ejecutar el algoritmo de Grover con menos iteraciones que las óptimas en condiciones ideales puede proporcionar una mayor probabilidad de éxito.}
\end{quote}

Se trata de una \textbf{hipótesis experimental}, no de una conclusión anticipada. Se evaluará mediante simulaciones con diferentes tamaños de problema, intensidades de ruido y números de iteraciones.

\section{Tecnologías y herramientas}

\subsection{Entorno de desarrollo}

\begin{center}
\begin{tabularx}{\textwidth}{@{}>{\bfseries}p{2.8cm}X@{}}
\toprule
Tecnología & Uso previsto \\
\midrule
Python & Implementación, automatización de experimentos y análisis de resultados. \\
Qiskit & Construcción de circuitos, oráculos y operadores de Grover. \\
Qiskit Aer & Simulación de circuitos ideales y ruidosos sin hardware cuántico físico. \\
\bottomrule
\end{tabularx}
\end{center}

Dentro de Qiskit Aer resultan especialmente relevantes los métodos de simulación \texttt{statevector}, para estudiar estados puros e ideales, y \texttt{density\_matrix}, que permite representar estados mixtos asociados a procesos ruidosos.

\subsection{Interfaces principales}

\begin{center}
\begin{tabularx}{\textwidth}{@{}>{\ttfamily}p{4.3cm}X@{}}
\toprule
\multicolumn{1}{@{}l}{\textbf{Herramienta}} & \textbf{Función} \\
\midrule
QuantumCircuit & Crear circuitos cuánticos. \\
grover\_operator() & Construir el operador de Grover a partir de un oráculo. \\
transpile() & Adaptar y compilar circuitos para un objetivo de ejecución. \\
AerSimulator & Ejecutar las simulaciones. \\
NoiseModel & Configurar modelos de ruido. \\
depolarizing\_error() & Crear errores despolarizantes. \\
ReadoutError & Configurar errores de lectura. \\
\bottomrule
\end{tabularx}
\end{center}

Para una implementación nueva, es preferible consultar la API actual de Qiskit y evitar copiar sin revisar tutoriales basados en componentes obsoletos.

\section{Recursos recomendados}

\begin{enumerate}
    \item \textbf{IBM Quantum Learning --- Grover's Algorithm.} Introducción a superposición, oráculos, difusión e iteraciones.\\
    \url{https://quantum.cloud.ibm.com/learning/en/modules/computer-science/grovers}
    \item \textbf{Qiskit --- Grover's Algorithm Tutorial.} Tutorial práctico para construir circuitos de Grover.\\
    \url{https://quantum.cloud.ibm.com/docs/en/tutorials/grovers-algorithm}
    \item \textbf{Qiskit Aer --- Building Noise Models.} Documentación para crear y aplicar modelos de ruido.\\
    \url{https://qiskit.github.io/qiskit-aer/tutorials/3_building_noise_models.html}
\end{enumerate}

Como referencia científica original puede consultarse:
\begin{quote}
L. K. Grover, ``A fast quantum mechanical algorithm for database search'', \emph{Proceedings of the 28th Annual ACM Symposium on Theory of Computing}, 1996. \url{https://arxiv.org/abs/quant-ph/9605043}
\end{quote}

\section{Conocimientos necesarios antes de implementar}

Antes de programar deberíamos poder explicar:
\begin{itemize}
    \item Qué problema resuelve Grover y en qué sentido mejora una búsqueda clásica.
    \item Cómo se forma una superposición con puertas Hadamard.
    \item Cómo funciona un oráculo de fase y por qué hace falta la difusión.
    \item Por qué existe un número de iteraciones adecuado y qué sucede si se supera.
    \item Cómo afectan los errores de puertas y los errores de lectura a la probabilidad de éxito.
    \item Para qué se utilizan Qiskit y Qiskit Aer.
\end{itemize}

\section{Siguiente paso propuesto}

Se recomienda comenzar implementando una búsqueda con \textbf{dos qubits} y el estado objetivo $\lvert 11\rangle$. Una vez comprobado su funcionamiento ideal, se podrá ampliar a tres y cuatro qubits. Posteriormente se incorporarán modelos de ruido para comparar probabilidades de éxito, tiempos de emulación y consumo de memoria.

\end{document}
